\section{Hashing de conjuntos}

\begin{frame}
  \vspace{1.5cm}
  \centering
  {\LARGE Hashing de conjuntos}
\end{frame}

\begin{frame}{XOR hashing}
  Conjuntos de \(\{0,\ldots,n-1\}\).
  Los \(r[i]\) son números aleatorios, elegidos una vez.
  \[
    \begin{aligned}
      h(\emptyset) &= 0 \\
      h(\{i\}) &= r[i] \\
      h(A \triangle B) &= h(A) \oplus h(B)
    \end{aligned}
  \]

  \vspace{-0.4em}
  \[
    \begin{aligned}
      h(\{a,b,c\})
        &= h(\{a\} \triangle \{b\} \triangle \{c\}) \\
        &= r[a] \oplus r[b] \oplus r[c]
    \end{aligned}
  \]

  Asociativa, conmutativa, con neutro \(0\).
  Cada conjunto es su inverso: \(h(A) \oplus h(A) = 0\).
\end{frame}

\begin{frame}{Insertar y borrar}
  Como \(r[i] \oplus r[i] = 0\), agregar y sacar son la misma cuenta.

  \medskip
  Agregar \(i\), sabiendo que no está:
  \[
    h \leftarrow h \oplus r[i]
  \]
  Sacar \(i\), sabiendo que está:
  \[
    h \leftarrow h \oplus r[i]
  \]

  \medskip
  Si \(i\) ya estaba y se agrega, el XOR lo saca.
  El hash no dice si \(i\) pertenece al conjunto.
\end{frame}
